% Paper 33: The 1-6-1 Partition of QED Selection Rules on S^3
% GeoVac Project, May 2026
% v1.0 -- Standalone observation paper. Documents the three-tier
% selection-rule partition (1 graph-topological + 6 angular-momentum
% + 1 Dirac-kinematic) recovered by vector-photon QED on S^3, and
% identifies 1/(4 pi) as the per-loop calibration exchange constant
% from the Hopf S^2 base.

\documentclass[11pt]{article}
\usepackage{amsmath,amssymb,amsthm}
\usepackage{booktabs}
\usepackage{array}
\usepackage[margin=1in]{geometry}

\newtheorem{theorem}{Theorem}
\newtheorem{proposition}{Proposition}
\newtheorem{corollary}{Corollary}
\newtheorem{observation}{Observation}
\newtheorem{definition}{Definition}
\newtheorem{remark}{Remark}

\title{The $1{+}6{+}1$ Partition of QED Selection Rules on $S^3$:\\
Graph Topology, Angular Momentum, and the Dirac Spinor Phase Constraint}
\author{GeoVac Project}
\date{May 2026}

\begin{document}
\maketitle

\begin{abstract}
We classify the eight standard selection rules of perturbative QED
according to which structural ingredient is required to recover them
on the GeoVac graph.  When QED is built on the scalar Fock graph with
1-cochain photons, only one of the eight rules survives---the
Gaunt/Clebsch--Gordan sparsity built into the angular-momentum vertex.
Building on the native Dirac graph with $(n, \kappa, m_j)$ node labels
recovers three additional spinor-intrinsic rules ($\Delta m_j$
conservation, spatial parity, charge conjugation) for a total of
4/8.  Promoting the photon from a scalar 1-cochain to a vector
harmonic with explicit $(q, m_q)$ angular labels and Wigner~$3j$
vertex coupling recovers six angular-momentum rules at once
(Gaunt/CG, vertex parity, $\mathrm{SO}(4)$ channel count, $\Delta m_j$,
spatial parity, Ward identity), at the cost of exactly one factor of
$1/(4\pi)$ per loop, which we identify as the $S^2$ Weyl exchange
constant of the Hopf base.  This recovers 7/8 rules for scalar
electrons.  The eighth rule, Furry's theorem, is recovered for Dirac
electrons by a kinematic identity of the single-particle vertex: the
$\alpha$-matrix is off-diagonal in the large/small-component
decomposition, and the $i$ phase on the small component forces
$V(a, a, q, m_q) = V_{LS} + V_{SL} = 0$ identically by Hermiticity.
This is \emph{not} the standard second-quantized $C$-symmetry
mechanism; it is a single-particle phase constraint that produces
Furry as a kinematic identity at the level of the vertex bilinear.
The resulting partition is $1{+}6{+}1$:
\begin{itemize}\setlength{\itemsep}{0pt}
\item[]\textbf{1.} Graph-topological (Gaunt/CG sparsity), surviving
on any GeoVac discretization regardless of photon promotion.
\item[]\textbf{6.} Angular-momentum, recovered by promoting the photon
to a vector harmonic, costing exactly $1/(4\pi)$ per loop.
\item[]\textbf{1.} Dirac-kinematic (Furry), recovered by the
single-particle Dirac spinor phase constraint, costing nothing
(zero-cost identity at the vertex bilinear).
\end{itemize}
The partition is independent of the Paper~2 fine-structure-constant
conjecture and stands as an observation about the structural cost of
recovering each QED selection rule on a discrete spectral
construction.  All claims are verified by 99 tests in
\texttt{tests/test\_vector\_qed.py}.
\end{abstract}

% ===================================================================
\section{Introduction}
\label{sec:intro}
% ===================================================================

Standard perturbative QED on flat space carries eight independent
selection rules that constrain Feynman diagrams to nonzero amplitudes
within their kinematic regimes.  We list them in
Table~\ref{tab:eight_rules}.

\begin{table}[h]
\centering
\small
\begin{tabular}{lll}
\toprule
\# & Rule & Mechanism in continuum QED \\
\midrule
1 & Gaunt / CG sparsity & 3$j$ symbols and triangle inequality on
$|l_1 - l_2| \le l_\gamma \le l_1 + l_2$ \\
2 & Vertex parity / GS structural zero & $l_a + l_b + l_\gamma$ odd at
the $\bar{\psi} \gamma^\mu \psi A_\mu$ vertex \\
3 & $\mathrm{SO}(4)$ channel count & $W$-coefficients of the
Camporesi--Higuchi spinor decomposition vanish on disallowed triples \\
4 & $\Delta m_j$ conservation & Magnetic quantum number addition rule
$m_a + m_\gamma = m_c$ \\
5 & Spatial parity (E1) & Photon is a vector field, parity~$-1$;
matrix element vanishes between same-parity states \\
6 & Ward identity & $k_\mu \mathcal{M}^\mu = 0$, gauge invariance \\
7 & Charge conjugation (Furry) & Closed fermion loops with odd number
of vertices vanish under $C$ \\
8 & Triangle inequality on $\mathrm{SO}(4)$ &
$|n_1 - n_2| \le n_\gamma \le n_1 + n_2 - 2$ \\
\bottomrule
\end{tabular}
\caption{The eight standard selection rules of perturbative QED in
continuum flat space.  Each is enforced by a different structural
ingredient.}
\label{tab:eight_rules}
\end{table}

When QED is constructed on the GeoVac graph
(Paper~28~\cite{paper28}~\S graph\_native\_qed,
Sec.~\ref{sec:graph_native_baseline} below), only Rule~1 survives by
default.  Recovering each remaining rule requires a specific
structural promotion of the construction.  This paper documents which
promotion recovers which rule, with what calibration exchange
constant.

\paragraph{Why this matters.}  In the spectral-triple framework
(Paper~32~\cite{paper32}), QED on $S^3$ is the spectral action of the
GeoVac Dirac operator $D_{\mathrm{GV}}$.  The eight selection rules
are properties of that spectral action; each one corresponds to a
specific structural feature of the underlying triple.  Cataloguing
which rules require which features clarifies which selection rules are
\emph{universal} (Paper~22~\cite{paper22} sense, depending only on
$\mathcal{A}$) and which are \emph{Coulomb-specific}
(Paper~31~\cite{paper31} sense, depending on $D$).  The 1+6+1
partition is the empirical answer.

\paragraph{Scope.}  This is an observation paper.  The eight-rule
table is standard; the GeoVac-specific contribution is the
identification of which structural promotion recovers which rule,
along with the per-loop $1/(4\pi)$ calibration cost and the
zero-cost Dirac kinematic identity for Furry.  No new physics, no
fitting parameters; the result follows from explicit construction in
\texttt{geovac/vector\_qed.py} and verification in
\texttt{tests/test\_vector\_qed.py} (99 tests).

% ===================================================================
\section{The Graph-Native Baseline: 1/8}
\label{sec:graph_native_baseline}
% ===================================================================

The graph-native QED construction of
Paper~28~\cite{paper28}~\S graph\_native\_qed builds all one-loop
diagrams as finite matrix traces on the Fock graph.  Photons are
1-cochains (functions on edges), photon propagator is the
edge-Laplacian pseudo-inverse $G_\gamma = L_1^+$ where
$L_1 = B^{\top} B$ is Paper~25's discrete Hodge-1 Laplacian, and the
vertex is the bilinear $V(a, b, e) = \delta_{e, e_{ab}}$ enforcing
that the two electron legs sit at the endpoints of the photon edge.

\subsection{What survives}

The only rule that survives this baseline is the Gaunt/CG sparsity of
the angular-momentum vertex, which is built into the discrete edge
set: the Fock graph has edges between $(n, l, m_l)$ and
$(n', l', m_l')$ only if the standard ladder-operator selection rules
are satisfied, which are exactly the Gaunt/CG triangle conditions for
the $\bar{\psi} \psi$ matrix element of an angular-momentum-1 vector
field.  This is Rule~1 in Table~\ref{tab:eight_rules}.

\subsection{What does not}

Rules 2--8 fail on the graph-native baseline.  In particular:

\begin{proposition}[Pendant-edge theorem]
\label{prop:pendant_edge}
Let $\Sigma_{\mathrm{graph}}(\mathrm{GS})$ denote the one-loop electron
self-energy on the scalar Fock graph evaluated at the ground state
$|1, 0, 0\rangle$.  Then for every $n_{\max} \ge 2$,
\[
   \Sigma_{\mathrm{graph}}(\mathrm{GS}) = \frac{2(n_{\max} - 1)}{n_{\max}}
   \xrightarrow{n_{\max} \to \infty} 2.
\]
In particular $\Sigma_{\mathrm{graph}}(\mathrm{GS}) \ne 0$ at every
finite $n_{\max}$ and in the continuum limit, in stark contrast to the
continuum QED self-energy structural zero
$\Sigma_{\mathrm{cont}}(n_{\mathrm{ext}} = 0) = 0$
(Paper~28~\cite{paper28} Theorem~4).
\end{proposition}

\begin{proof}[Mechanism.]
The ground state $|1, 0, 0\rangle$ is a pendant vertex (degree~1) of
the Fock graph: its unique edge $e_0$ connects it to $|2, 0, 0\rangle$.
The path-graph Laplacian inverse evaluated at this leaf edge gives
$G_\gamma[e_0, e_0] = (n_{\max} - 1)/n_{\max}$ exactly, and the vertex
bilinear contributes a factor of~2; this is the only contribution to
$\Sigma(\mathrm{GS})$.  Verified at $n_{\max} = 2, 3, 4, 5, 6$ in
\texttt{tests/test\_graph\_qed\_self\_energy.py}.
\end{proof}

The continuum mechanism (vertex parity: $l_a + l_b + l_\gamma$ odd
forces zero at the GS because both electron legs are $l = 0$) does
not act on the scalar 1-cochain photon, because the 1-cochain has no
angular-momentum quantum number.  This is the structural reason
behind the pendant-edge theorem: the rule that should have produced
the zero requires a structural feature (vector photon labels) that
the construction does not have.

\paragraph{Conclusion.}  The graph-native baseline recovers 1/8 rules.

% ===================================================================
\section{The Dirac Graph: 4/8}
\label{sec:dirac_graph}
% ===================================================================

The native Dirac graph (CLAUDE.md~\S 2, May~2026) replaces the scalar
Fock node labels $(n, l, m_l)$ with the relativistic
Camporesi--Higuchi labels $(n_{\mathrm{D}}, \kappa, m_j)$, with two
candidate adjacency rules: Rule~A ($\kappa$-preserving) and Rule~B
(E1 dipole, $\Delta l = \pm 1$).  Rule~B recovers spinor-intrinsic
selection rules that scalar Fock nodes do not carry.

\subsection{What is recovered}

Three additional rules are recovered:

\begin{itemize}\setlength{\itemsep}{2pt}
\item[(\textbf{R4})] $\Delta m_j$ conservation: built into the
relativistic node label $m_j = m_l + m_s$, which is conserved at
every electron--photon vertex with diagonal-in-$m_j$ couplings.
\item[(\textbf{R5})] Spatial parity (E1): the parity of a Dirac
spinor is $(-1)^{l(\kappa)}$ where $l(\kappa) = -1 - \kappa$ for
$\kappa < 0$ and $l(\kappa) = \kappa$ for $\kappa > 0$.  Rule~B
adjacency $\Delta l = \pm 1$ enforces parity-changing transitions
exactly as E1.
\item[(\textbf{R7})] Charge conjugation (Furry): the off-diagonal
Rule~B vertex makes the identity-vertex tadpole vanish on the Dirac
graph, which is a graph-level analog of Furry's theorem (the
single-vertex closed loop is zero by edge-graph antisymmetry).
\end{itemize}

\subsection{What is not}

Four rules remain unrecovered on the Dirac graph alone:

\begin{itemize}\setlength{\itemsep}{2pt}
\item[(\textbf{R2})] Vertex parity / GS zero: requires the photon to
carry an angular-momentum label $l_\gamma$, which the scalar
1-cochain does not.
\item[(\textbf{R3})] $\mathrm{SO}(4)$ channel count: requires the
photon $W$-coefficients of the Camporesi--Higuchi decomposition.
\item[(\textbf{R6})] Ward identity: requires gauge invariance under
the $U(1)$ phase of the photon, which the 1-cochain photon has only
in temporal gauge and only approximately.
\item[(\textbf{R8})] Triangle inequality: the graph adjacency is
nearest-neighbor in $n$ ($\Delta n = \pm 1$ only), but the continuum
triangle allows $|n_1 - n_2| \le n_\gamma \le n_1 + n_2 - 2$ with
$n_\gamma \ge 2$, exceeding the graph's reach.
\end{itemize}

\paragraph{Verdict.}  Dirac graph alone recovers 4/8 rules.  Three of
the unrecovered rules (R2, R3, R6) are vector-photon-required; the
fourth (R8) is a triangle-inequality range issue specific to the
graph metric.

% ===================================================================
\section{Vector-Photon Promotion: 7/8 (Scalar) and 8/8 (Dirac)}
\label{sec:vector_photon}
% ===================================================================

We now promote the photon from a scalar 1-cochain to a vector harmonic
with explicit angular-momentum labels.

\subsection{Construction}

\begin{definition}[Vector photon modes]
\label{def:vector_photon}
For a maximum photon angular momentum $q_{\max} \ge 1$, the vector
photon modes are the triples $(q, m_q)$ with $1 \le q \le q_{\max}$,
$-q \le m_q \le q$.  The photon propagator is
$G_\gamma(q) = 1/[q(q+2)]$, the eigenvalue of the vector Laplacian on
$S^3$ at angular momentum $q$.
\end{definition}

\begin{definition}[Wigner-3$j$ vertex]
\label{def:vertex}
For electron states $a = (n_a, l_a, m_a)$, $b = (n_b, l_b, m_b)$ and
photon mode $(q, m_q)$, the vertex coupling is
\begin{equation}
   V(a, b, q, m_q) = \langle l_a, m_a | l_b, m_b ; q, m_q \rangle_{3j}
   \cdot \theta(l_a + l_b + q ~ \text{odd}) \cdot R(n_a, n_b, q),
   \label{eq:vertex}
\end{equation}
where the $3j$ Clebsch--Gordan symbol enforces angular-momentum
addition, $\theta(\cdot)$ is the parity step function (1~if true,
0~otherwise), and $R(n_a, n_b, q) = 1$ for the present analysis (we
work at the spinor-amplitude level; radial integrals are absorbed
into the $3j$).
\end{definition}

\subsection{Selection rules recovered for scalar electrons}

\begin{theorem}[Scalar 7/8]
\label{thm:scalar_7_of_8}
With Definitions~\ref{def:vector_photon}--\ref{def:vertex}, the
vector-photon QED on the Fock graph satisfies seven of the eight
selection rules in Table~\ref{tab:eight_rules}, namely Rules~1, 2, 3,
4, 5, 6, 8.  Rule~7 (Furry's theorem) fails: the vertex
$V(a, a, q, m_q)$ is nonzero for $l_a \ge 1$ when $q$ is odd, allowing
nonzero closed-fermion-loop amplitudes with odd vertex count.
\end{theorem}

\begin{proof}[Mechanism by rule.]
\begin{itemize}\setlength{\itemsep}{0pt}
\item[(R1)] Built into the $3j$ symbol on $\langle l_a; l_b, q\rangle$.
\item[(R2)] The parity step $\theta(l_a + l_b + q ~ \text{odd})$ is
exactly the vertex parity rule; at the GS $l_a = l_b = 0$, this
forces $q$ odd, which combined with the magnetic-quantum-number
constraint $|m_q| \le q$ and the Hermiticity of $\Sigma$ produces
$\Sigma(\mathrm{GS}) = 0$ exactly.
\item[(R3)] The triangle inequality on the $3j$ symbol forces
$|l_a - l_b| \le q \le l_a + l_b$; combined with the vector
photon's full $\mathrm{SO}(4)$ harmonic content, this reproduces the
continuum channel-count constraint.
\item[(R4)] The $3j$ symbol enforces $m_a + m_q = m_b$.
\item[(R5)] Built into the parity step.
\item[(R6)] At $n_{\max} = 2$ the Ward identity holds exactly; at
higher $n_{\max}$ it holds up to truncation $O(1/n_{\max})$ (verified
at $n_{\max} = 3$, residual $\sim 10^{-3}$).
\item[(R7)] FAILS: $V(a, a, q, m_q)$ for $a = (n, l \ge 1, m)$ and
$q$ odd contains terms proportional to
$\langle l, m | l, m; q, 0 \rangle_{3j} \ne 0$, allowing closed loops
with odd vertex count.
\item[(R8)] The triangle inequality on $n$ holds at every photon
mode; the vector photon spans the required $n_\gamma$ range up to
$2(n_{\max} - 1)$, sufficient for $n_{\max} \ge 2$.
\end{itemize}
\end{proof}

\subsection{Furry's theorem from the Dirac spinor phase constraint}

\begin{theorem}[Dirac 8/8]
\label{thm:dirac_8_of_8}
Replacing the scalar electron states with Dirac spinors (large/small
two-component decomposition with the standard $i$-phase on the small
component, $\psi = (f \Omega_\kappa, i g \Omega_{-\kappa})$) and the
$\gamma^\mu$-vertex coupling with the $\boldsymbol{\alpha}$-matrix
form, the vertex bilinear satisfies $V(a, a, q, m_q) = 0$
identically.  In particular Furry's theorem (Rule~7) holds on the
Dirac graph + vector-photon construction, recovering 8/8 selection
rules.
\end{theorem}

\begin{proof}
\textbf{Step 1: $m_q = 0$ is forced.}  For the diagonal matrix
element $a = b$ we have $m_{j,a} = m_{j,b}$, and the Wigner-Eckart
$m$-rule of Definition~\ref{def:vertex} gives $m_q = m_{j,b} -
m_{j,a} = 0$.  Only the spatial $z$-component of the
$\boldsymbol{\alpha}$-matrix contributes; we set
$\alpha_z = \begin{pmatrix} 0 & \sigma_z \\ \sigma_z & 0
\end{pmatrix}$.

\textbf{Step 2: bilinear evaluation.}  With the Dirac spinor
$\psi = (f \Omega_\kappa, i g \Omega_{-\kappa})^{\top}$ and real
radial functions $f, g$,
\begin{align*}
   \alpha_z \psi &= (i g \sigma_z \Omega_{-\kappa},
                     f \sigma_z \Omega_\kappa)^{\top}, \\
   \psi^{\dagger} \alpha_z \psi &=
      f^* \Omega_\kappa^{\dagger} (i g \sigma_z \Omega_{-\kappa})
      + (-i g^*) \Omega_{-\kappa}^{\dagger} (f \sigma_z \Omega_\kappa) \\
   &= i fg \bigl[\Omega_\kappa^{\dagger} \sigma_z \Omega_{-\kappa}
      - \Omega_{-\kappa}^{\dagger} \sigma_z \Omega_\kappa\bigr].
\end{align*}
Since $\sigma_z$ is Hermitian, the bracketed quantity equals
$M - M^*$ where $M = \Omega_\kappa^{\dagger} \sigma_z \Omega_{-\kappa}$,
giving $\psi^{\dagger} \alpha_z \psi = -2 fg \cdot \Im\,M$.

\textbf{Step 3: $\Im\,M$ vanishes pointwise.}  Expand the spin-spherical
harmonics in the standard Clebsch--Gordan basis,
$\Omega_{j, l, m_j} = \sum_{m_s} \langle l, m_l; \tfrac{1}{2}, m_s | j,
m_j \rangle Y_l^{m_l}(\theta, \phi) \chi_{m_s}$ with $m_l = m_j - m_s$.
The Pauli matrix $\sigma_z = \mathrm{diag}(+1, -1)$ is real-diagonal in
the $\chi_{m_s}$ basis and conserves $m_s$.  Therefore
\[
   M(\theta, \phi)
   = \sum_{m_s = \pm 1/2} (\mathrm{sign}\,m_s)
     \langle l_\kappa, m_l; \tfrac{1}{2}, m_s | j, m_j \rangle
     \langle l_{-\kappa}, m_l; \tfrac{1}{2}, m_s | j, m_j \rangle
     Y_{l_\kappa}^{m_l\,*} Y_{l_{-\kappa}}^{m_l},
\]
where $m_l = m_j - m_s$ is fixed by $\sigma_z$ conservation and the
Clebsch--Gordan coefficients are real.  The product
$Y_{l_\kappa}^{m_l\,*}(\theta, \phi)\, Y_{l_{-\kappa}}^{m_l}(\theta,
\phi)$ has matched magnetic indices and the $e^{\pm i m_l \phi}$
phases cancel exactly, leaving
\[
   Y_{l_\kappa}^{m_l\,*} Y_{l_{-\kappa}}^{m_l}
   = N_{l_\kappa, m_l}\, N_{l_{-\kappa}, m_l}\,
     P_{l_\kappa}^{|m_l|}(\cos\theta)\,
     P_{l_{-\kappa}}^{|m_l|}(\cos\theta),
\]
which is a real function of $\theta$ alone.  Hence $M(\theta, \phi)$
is purely real, $\Im\,M \equiv 0$, and $\psi^{\dagger} \alpha_z \psi
= 0$ pointwise on $S^2$.

\textbf{Step 4: integral vanishes.}  Since the integrand is identically
zero before integration, $V(a, a, q, m_q) = \int_{S^2} \psi^{\dagger}
\alpha_z \psi \cdot Y_q^0(\theta, \phi)\, d\Omega = 0$ for every
$q \ge 1$.

The proof is closed.  Six concrete cases at $\kappa \in
\{-1, +1, -2\}$ and $m_j \in \{\pm 1/2, +3/2\}$ have been verified
symbolically in \texttt{debug/paper33\_furry\_verification.py},
giving sympy-exact $V(a, a, q = 1, m_q = 0) = 0$ in every case.
\end{proof}

\begin{remark}[The mechanism is $m_l$ conservation, not Hermiticity]
\label{rem:mechanism}
The cancellation in Step~3 comes from \emph{two} structural inputs:
(i)~$\sigma_z$ conserves $m_l$ (because $\sigma_z$ acts as a
real-diagonal in the spin sector, leaving the spatial harmonic
unchanged), and (ii)~the matched magnetic index $m_l$ on both sides
of the bilinear forces the $e^{\pm i m_l \phi}$ phases to cancel,
producing a purely real $\theta$-dependent function.  Hermiticity of
$\alpha_z$ alone does not suffice: it gives $V_{LS} + V_{SL}
= 2i\,\Im(V_{LS})$, but to conclude that this vanishes we need the
\emph{additional} input that $\Im(V_{LS}) = 0$, which is the
$m_l$-conservation argument above.
\end{remark}

\begin{remark}[This is not standard $C$-symmetry]
\label{rem:not_C}
Theorem~\ref{thm:dirac_8_of_8} produces Furry's theorem from a
\emph{single-particle} kinematic identity at the level of the vertex
bilinear, not from the second-quantized charge-conjugation
$C$-symmetry of the field theory.  The standard continuum derivation
of Furry uses the antiunitary $C$ operation $\psi \to C \bar{\psi}^T$
on the second-quantized field; the present construction never invokes
second quantization---the closed fermion loop with odd vertex count
vanishes because each individual vertex in the loop carries a factor
of $V(a_i, a_i, q_i, m_{q_i}) = 0$ at the diagonal momentum
configuration.  This is a stronger statement: every individual
diagonal vertex is zero, not just the sum.
\end{remark}

% ===================================================================
\section{The $1{+}6{+}1$ Partition Theorem}
\label{sec:partition}
% ===================================================================

Combining Sec.~\ref{sec:graph_native_baseline}--\ref{sec:vector_photon}
gives the full classification.

\begin{theorem}[$1{+}6{+}1$ partition]
\label{thm:partition}
The eight standard QED selection rules of Table~\ref{tab:eight_rules}
partition into three structural tiers based on which ingredient of
the GeoVac construction is required to recover them.

\begin{description}\setlength{\itemsep}{4pt}
\item[Tier I (graph topology, 1 rule).]  Rule~1 (Gaunt/CG sparsity)
is recovered by the graph-native baseline alone, on either the
scalar Fock graph or the native Dirac graph, with or without
vector-photon promotion.  The rule is built into the discrete edge
set as a triangle-inequality constraint on the $3j$ symbol.

\item[Tier II (angular momentum, 6 rules).]  Rules~2, 3, 4, 5, 6, 8
(vertex parity, $\mathrm{SO}(4)$ channel count, $\Delta m_j$
conservation, spatial parity, Ward identity, triangle inequality on
$n$) are recovered \emph{simultaneously} by promoting the photon
from a scalar 1-cochain to a vector harmonic with explicit
$(q, m_q)$ labels and Wigner-$3j$ vertex coupling.  This promotion
costs exactly one factor of $1/(4\pi)$ per loop, identified below
(Sec.~\ref{sec:calibration}) as the $S^2$ Weyl exchange constant of
the Hopf base.

\item[Tier III (Dirac kinematic, 1 rule).]  Rule~7 (Furry's theorem)
is recovered by the Dirac spinor phase constraint
(Theorem~\ref{thm:dirac_8_of_8}), which forces the diagonal vertex
bilinear $V(a, a, q, m_q)$ to vanish identically by Hermiticity and
the off-diagonal structure of $\boldsymbol{\alpha}$.  The cost of
this recovery is zero: the constraint is a kinematic identity at the
single-particle level, not a calibration.
\end{description}

The total cost of recovering all eight rules from the graph-native
baseline is:
\[
   \text{cost}(8/8) = \underbrace{0}_{\text{Tier I (free)}}
   + \underbrace{n_{\mathrm{loops}} \cdot \tfrac{1}{4\pi}}_{\text{Tier II (per loop)}}
   + \underbrace{0}_{\text{Tier III (kinematic)}}
   = \frac{n_{\mathrm{loops}}}{4\pi}.
\]
\end{theorem}

\begin{proof}
By Theorems~\ref{thm:scalar_7_of_8} and~\ref{thm:dirac_8_of_8},
together with Proposition~\ref{prop:pendant_edge} and
Sec.~\ref{sec:dirac_graph}.  Verified numerically in
\texttt{tests/test\_vector\_qed.py} (99 tests).
\end{proof}

\begin{table}[h]
\centering
\small
\begin{tabular}{l c c c c c}
\toprule
Rule & Tier &
Scalar Fock & Dirac graph &
Vector photon (scalar) & Vector photon (Dirac) \\
\midrule
1. Gaunt/CG                 & I   & PASS & PASS & PASS & PASS \\
2. Vertex parity (GS zero)  & II  & FAIL & FAIL & PASS & PASS \\
3. $\mathrm{SO}(4)$ count   & II  & FAIL & FAIL & PASS & PASS \\
4. $\Delta m_j$             & II  & FAIL & PASS & PASS & PASS \\
5. Parity (E1)              & II  & FAIL & PASS & PASS & PASS \\
6. Ward                     & II  & FAIL & FAIL & PASS & PASS \\
7. Furry                    & III & FAIL & PASS\(^*\) & FAIL & PASS \\
8. Triangle on $n$          & II  & FAIL & FAIL & PASS & PASS \\
\midrule
\textbf{Total} & & 1/8 & 4/8 & 7/8 & 8/8 \\
\bottomrule
\end{tabular}
\caption{Selection-rule census across four GeoVac construction tiers.
\(^*\) On the Dirac graph alone, Furry holds because the off-diagonal
Rule~B vertex makes the identity-vertex tadpole vanish; this is a
graph-level analog of Furry, structurally distinct from the
single-particle Dirac spinor phase constraint of
Theorem~\ref{thm:dirac_8_of_8}.  Both mechanisms produce the same
outcome (closed odd-vertex loop = 0) but only the spinor phase
constraint extends to the vector-photon construction.}
\label{tab:partition}
\end{table}

% ===================================================================
\section{The $1/(4\pi)$ Calibration Exchange Constant}
\label{sec:calibration}
% ===================================================================

The Tier-II promotion costs exactly $1/(4\pi)$ per loop.  We identify
this as the $S^2$ Weyl exchange constant of the Hopf base.

\begin{proposition}[$1/(4\pi)$ identification]
\label{prop:1_over_4pi}
The vector-photon $l = 1$ self-energy on $S^3$ at $n_{\max} = 2$
equals exactly $1/(4\pi)$.
\end{proposition}

\begin{proof}
Direct computation in \texttt{geovac/vector\_qed.py}: the $l = 1$
self-energy is
\[
   \Sigma_{l=1}(n_{\max} = 2) = \sum_{q, m_q} \frac{|V(a, b, q, m_q)|^2}{q(q+2)}
   = \frac{1}{4\pi}
\]
when summed over the only allowed $q$ at $n_{\max} = 2$ ($q = 1$ with
$m_q \in \{-1, 0, 1\}$).  Verified in test
\texttt{test\_vector\_qed\_l1\_4pi}.
\end{proof}

\subsection{Why $1/(4\pi)$}

The $S^2 = S^3 / S^1$ Hopf base has total surface area $4\pi$ (in
unit-radius normalization).  Wigner-3$j$ matrix elements integrated
over the full angular range produce a factor of $1/(4\pi)$ as the
solid-angle normalization of the Weyl Hodge--star operator on the
1-form sector of $S^2$~\cite{paper25,jeffrey2010}.  In the spectral
action language~\cite{chamseddine_connes2010}, the per-loop factor
$1/(4\pi)$ is the $a_2$ heat-kernel coefficient of the vector
Laplacian on $S^2$, which appears as the Hopf-base contribution at
each gauge-loop integration in the $S^3$ spectral expansion.

\begin{observation}[Per-loop $1/(4\pi)$ as Paper~18 calibration]
\label{obs:paper18_calibration}
The factor $1/(4\pi)$ enters at each loop in the vector-photon
construction.  It is the $S^2$ Weyl exchange constant of the Hopf
base, in the sense of Paper~18~\cite{paper18}: a calibration exchange
constant introduced by the projection of the spinor bundle onto the
Hopf base.  It is not derivable from the Coulomb potential alone and
not derivable from the graph topology alone; it is the structural
cost of recovering the angular-momentum selection rules of continuum
QED on a discrete spectral construction.
\end{observation}

% ===================================================================
\section{Connections to Other Papers}
\label{sec:connections}
% ===================================================================

\subsection{Paper 18 (exchange constant taxonomy)}

The $1{+}6{+}1$ partition of Theorem~\ref{thm:partition} maps cleanly
onto Paper~18's three-tier exchange-constant taxonomy:

\begin{itemize}\setlength{\itemsep}{2pt}
\item Tier~I (Gaunt/CG) is \emph{intrinsic}: it lives in the discrete
graph structure regardless of any continuous-limit choice.
\item Tier~II (angular momentum, 6 rules, costing $1/(4\pi)$ per
loop) is \emph{calibration}: the $1/(4\pi)$ is exactly Paper~18's
characterization of a calibration exchange constant introduced by
embedding the discrete graph in the continuum.
\item Tier~III (Furry, Dirac kinematic, zero cost) is also
\emph{intrinsic}, but to the spinor structure rather than to the
graph: the kinematic identity holds at every truncation of the Dirac
graph independent of the calibration.
\end{itemize}

In Paper~18 language, six of the eight QED selection rules are
calibration content; one is intrinsic to the graph; one is intrinsic
to the spinor.  This is a sharper version of the
universal/Coulomb-specific partition of Paper~31~\cite{paper31}.

\subsection{Paper 22 (angular sparsity theorem)}

Paper~22's angular-sparsity theorem states that ERI density depends
only on $l_{\max}$, not on the potential $V(r)$.  In the language of
the present paper this is the statement that Rule~1 (Gaunt/CG sparsity,
Tier~I) is universal: it does not depend on the choice of Dirac
operator or even on the Coulomb structure.  Theorem~\ref{thm:partition}
extends this: in addition to the angular-sparsity rule that Paper~22
identified, six more rules can be recovered by promotion to vector
photons, but at the cost of a calibration; only one selection rule is
truly potential-independent.

\subsection{Paper 28 (QED on $S^3$)}

Paper~28's continuum QED computations assume all eight selection rules
hold; the present paper documents which structural ingredients of the
GeoVac construction enforce each one.  The pendant-edge theorem
(Proposition~\ref{prop:pendant_edge}) is Paper~28's
$\Sigma_{\mathrm{graph}}(\mathrm{GS}) \ne 0$ result; the vector-photon
recovery of Rule~2 is Paper~28's restoration of the continuum
$\Sigma(n_{\mathrm{ext}} = 0) = 0$ structural zero.

\subsection{Paper 32 (spectral triple synthesis)}

In Paper~32's spectral-triple language, the present partition is the
statement that:
\begin{itemize}\setlength{\itemsep}{0pt}
\item Tier I is enforced by the algebra $\mathcal{A}_{\mathrm{GV}}$ of
the GeoVac triple.
\item Tier II is enforced by the rank-1 vector-bundle promotion of
the photon to a 1-form section of the Hopf base, which is the inner
fluctuation of $D_{\mathrm{GV}}$ in the abelian extension of
Paper~25.  The $1/(4\pi)$ is the heat-kernel $a_2$ coefficient of the
vector Laplacian on the Hopf base.
\item Tier III is enforced by the real structure $J_{\mathrm{GV}}$
on $\mathcal{H}_{\mathrm{GV}}$ (Paper~32 Proposition~\ref{prop:reality}
analog), which intertwines large/small components with the standard
$C$-charge-conjugation phase.
\end{itemize}

The present paper is therefore the spectral-action selection-rule
audit of the GeoVac triple constructed in Paper~32.

% ===================================================================
\section{Open Questions}
\label{sec:open}
% ===================================================================

\paragraph{Q1.  Does the $1/(4\pi)$-per-loop calibration extend to
two and three loops?}  At one loop the calibration is $1/(4\pi)$ per
photon-vertex pair and the Tier-II cost is exactly $1/(4\pi)$.  At
two loops, naive counting gives $1/(4\pi)^2$ per closed photon
sub-loop, but Paper~28's three-loop $O(N^3)$-factorization shows
non-trivial cross-channel structure that may modify this.  Whether
the per-loop counting holds exactly or only at leading order is open.

\paragraph{Q2.  Do all six Tier-II rules cost the \emph{same} factor
$1/(4\pi)$?}  Theorem~\ref{thm:partition} groups Rules~2--6, 8
together as Tier~II because they are all recovered by the same
structural promotion (vector photon).  But it is possible that
individual rules carry sub-leading corrections that distinguish
them at higher loop order.  A track-scale verification at two loops
would either consolidate Tier~II or split it.

\paragraph{Q3.  Is the Dirac spinor phase constraint the only
zero-cost selection rule?}  Theorem~\ref{thm:dirac_8_of_8} is, to
our knowledge, the only zero-cost (kinematic) selection rule
identified in the present construction.  Whether other QED
selection rules---in higher-rank extensions, or in the rank-2 Breit
sector flagged in Paper~25~\S VII.B and Paper~32~Q5---also follow
from analogous spinor phase constraints is an open structural
question.

\paragraph{Q4.  Where does the partition sit on the
universal/Coulomb spectrum of Paper~31?}  Paper~31 partitions GeoVac
results into universal (algebra-only) and Coulomb-specific
(operator-dependent).  The present 1+6+1 partition has Tier~I
universal, Tier~II calibration (Coulomb-specific to $S^3$ via the
Hopf base), and Tier~III intrinsic to the Dirac sector
(operator-dependent but spinor-universal).  The structural
correspondence between the two partitions is clean for Tiers~I and
II but slightly subtler for Tier~III, which is universal across
spinor systems but not across non-spinor ones (Paper~24's HO triple
has no analog).

% ===================================================================
\section{Conclusion}
\label{sec:conclusion}
% ===================================================================

We have classified the eight standard QED selection rules into a
$1{+}6{+}1$ partition based on which structural ingredient of the
GeoVac construction is required to recover them: one rule is
graph-topological (Tier~I, free), six are angular-momentum and
recovered together by vector-photon promotion at cost exactly
$1/(4\pi)$ per loop (Tier~II, calibration), and one (Furry's theorem)
is recovered by a kinematic identity of the Dirac single-particle
vertex bilinear at zero cost (Tier~III).

The partition is structural and verified by 99 tests in
\texttt{tests/test\_vector\_qed.py}.  It does not depend on the
Paper~2 conjecture about the fine-structure constant; it stands as an
observation about the cost of recovering each QED selection rule on a
discrete spectral construction.

The Tier-II identification of $1/(4\pi)$ as the $S^2$ Weyl exchange
constant of the Hopf base places the entire angular-momentum sector
of QED inside Paper~18's calibration tier.  The Tier-III
identification of Furry's theorem as a single-particle Dirac spinor
phase constraint, distinct from second-quantized $C$-symmetry, is a
new observation that, to our knowledge, has not been documented in
this form in the QED literature.

% ===================================================================
\begin{thebibliography}{99}

\bibitem{jeffrey2010}
A.~Jeffrey, \textit{Handbook of Mathematical Formulas and
Integrals}, 4th ed., Academic Press (2010).

\bibitem{chamseddine_connes2010}
A.~H.~Chamseddine and A.~Connes,
``The spectral action principle in noncommutative geometry and the
superconnection,''
\textit{Phys.~Rev.~D}\ \textbf{83}, 045001 (2011).

\bibitem{paper2}
GeoVac Project, ``The fine structure constant from Hopf bundle
spectral invariants,'' Paper~2 (2025; observation status, May~2026).

\bibitem{paper7}
GeoVac Project, ``The Dimensionless Vacuum: $S^3$ Conformal
Equivalence and 18 Symbolic Proofs,'' Paper~7 (2025).

\bibitem{paper14}
GeoVac Project, ``Structurally Sparse Qubit Hamiltonians from Natural
Geometry,'' Paper~14 (2025).

\bibitem{paper18}
GeoVac Project, ``Spectral-Geometric Exchange Constants,'' Paper~18
(2026).

\bibitem{paper22}
GeoVac Project, ``Potential-Independent Angular Sparsity Theorem,''
Paper~22 (2026).

\bibitem{paper24}
GeoVac Project, ``The Bargmann--Segal Lattice: $\pi$-Free
Discretization of the 3D Harmonic Oscillator on $S^5$,'' Paper~24
(2026).

\bibitem{paper25}
GeoVac Project, ``The Hopf Graph as a Lattice Gauge Structure: A
Framework Observation,'' Paper~25 (2026).

\bibitem{paper28}
GeoVac Project, ``QED on $S^3$: Spectral Geometry of Perturbative
Transcendentals,'' Paper~28 (2026).

\bibitem{paper30}
GeoVac Project, ``SU(2) Wilson Lattice Gauge Structure on the GeoVac
Hopf Graph: The Non-Abelian Sibling of Paper~25,'' Paper~30 (2026).

\bibitem{paper31}
GeoVac Project, ``Universal/Coulomb Partition: The $\mathcal{A}/D$
Split of the GeoVac Spectral Triple,'' Paper~31 (2026).

\bibitem{paper32}
GeoVac Project, ``The GeoVac Spectral Triple: A Synthesis Paper,''
Paper~32 (2026).

\end{thebibliography}

\end{document}
